% *****************************************************************************
% Copyright 2025 ULFPP

% Author: A.S.Grm
% Description: This is an official LaTeX ULFPP thesis disposition template.
% Version: 2.0
% Date: 25/11/2025
%
% ----------------------------------------------------------------------
%                             !!! POMEMBNO !!!
%       
%   Predloga je narejena za prevajanje s poljubnim prevajalnikom:
%        - pdflatex, xelatex, luatex
%   kjer je potrebno uporabiti biber ukaz za izdelavo bibliografije !!!
% ----------------------------------------------------------------------
%
%                *** Izdelava PDF/A dokument ***
%
%
% pdfLatex generira navaden PDF, ki ga je potrebno pretvorit v PFD/A
%
% pretvorba v PDF/A online: 
%           https://www.pdfforge.org/online/en/pdf-to-pdfa
% 
% validacija PDF/A: 
%           https://www.pdfforge.org/online/en/validate-pdfa
%
% *****************************************************************************


% *****************************************************************************
%
%              !!! IZBIRA USTREZNEGA PROGRAMA in/ali SMERI !!!
%
%  Najprej je potrebno določit deklaracijo študija v 
%
%    \documentclass[<program študija>,<smer>]{fppdisp}
%
% - zamenjava za <program študija> z ustreznim programom
% - če je potrebna smer pa še vpis smeri: SMER={splošna,vojaška}
%
%  primer: \documentclass[pomnav,SMER=splošna]{fppdisp} 
%
%  - pomnav:   pomorstvo vsš, 1. stopnja, navtika 
%      potrebno dodati še SMER={splošna ali vojaška} 
%   
%  - pomstr:   pomorstvo vsš, 1. stopnja, pomorsko strojništvo
%      potrebno dodati še SMER={splošna ali vojaška}  
%   
%  - pom2:     pomorstvo, 2. stopnja
%      izbriši opcijo SMER v \setkeys{fppdisp}
%
%  - promtpl:  promet uni, 1. stopnja, TPL
%      izbriši opcijo SMER v \setkeys{fppdisp}
%
%  - prompttl: promet vsš, 1. stopnja, PTTL
%      potrebno dodati še SMER={splošna ali vojaška}
%
%  - prom2:    promet, 2. stopnja
%      izbriši opcijo SMER v \setkeys{fppdisp}
%
% *** Pomorstvo ***
\documentclass[pomnav,SMER=splošna]{fppdisp}
%\documentclass[pomstr,SMER=splošna]{fppdisp}
%\documentclass[pom2]{fppdisp}
% *** Promet ***
%\documentclass[prompttl,SMER=splošna]{fppdisp}
%\documentclass[promtpl]{fppdisp}
%\documentclass[prom2]{fppdisp}

% Zamenjajte z vašimi podatki, vsi podatki so potrebni!
\setkeys{fppdisp}{
	NASLOV={Preverjanje točnosti simulacij plovbe},
	AVTOR={Janez NOVAK},
	STEVILKA=123456,
	MENTOR={doc.dr.~Miha BONČA},
	MESEC=November,
	LETO=2025
}


% *****************************************************************************
% DODATNI PAKETI in FUNKCIONALNOSTI
% 
% WARNING: minimalno, da v OverLeaf prevede v kratkem času. Če boste dodali
%          preveč paketov, OverLeaf ne bo pravočasno prevedel!      
% 
% *****************************************************************************
\usepackage{enumitem}   % okolje za naštevanje
\usepackage{siunitx}    % dodatno okolje za pravilno pisanje števil z enoto
\usepackage{csquotes}   % biber potrebuje

% *** Caption to be bold face ***
\usepackage[labelfont=bf]{caption}

% *** Serif font za poglavja ***
\usepackage{titlesec}
\titleformat{\section}
{\normalfont\sffamily\Large\bfseries}
{\thesection}{1em}{}
\titleformat{\subsection}
{\normalfont\sffamily\large\bfseries}
{\thesubsection}{1em}{}


% *****************************************************************************
% LITERATURA 
% *****************************************************************************
\usepackage{csquotes}   % biber potrebuje
\usepackage[style=numeric,backend=biber,sorting=none]{biblatex}

% potrebno za deljenje dolgih URL naslovov
\setcounter{biburllcpenalty}{7000}
\setcounter{biburlucpenalty}{8000}

% ime datoteke, ki vsebuje navedeno literaturo
\addbibresource{literatura.bib}


% *****************************************************************************
% ZAČETEK VSEBINE
% *****************************************************************************
\begin{document}

% *********************
% *** Novo poglavje ***
% *********************
\section{Uvod -- Opis problema}

Uvodni del naloge ponavadi vsebuje:

\begin{itemize}[nosep]
	\item splošen pregled področja naloge,
	\item osnovno literaturo področja naloge,
	\item kaj bomo v nalogi naredili,
	\item kako je naloga sestavljena.
\end{itemize}

Namen uvodnega dela je opisati ozadje diplomskega dela \cite{froude1888resistance}, kaj so glavna izhodišča in hipoteza \cite{kalisnik}. Nadaljujete s splošnim opisom pristopa reševanja naloge, kjer se sklicujete na vire \cite{web:lang:stats}, ki opisujejo določene postopke \cite{glob, kalisnik}, rezultate, algoritme in podobno \cite{wiki, zbornik, vec-avtorjev}, iz katerih se boste navdihovali oziroma nadaljevali vaše raziskovalno delo v nalogi \cite{lang}.


% *********************
% *** Novo poglavje ***
% *********************
\section{Cilji in namen naloge}

V tem poglavju natančno opišete kaj so cilji naloge.


% *********************
% *** Novo poglavje ***
% *********************
\section{Predvidene metode dela}

V tem poglavju opišete katere metode boste uporabili pri delu, kako in zakaj jih boste uporabili. Seznam metod je sledeči:

\begin{itemize}
	\item \textbf{kvalitativne metode}
	\begin{itemize}
		\item Kaj so?\\
		Metode, ki se osredotočajo na neštevilske podatke (besedilo, opisi, izkušnje, mnenja).
		\item Namen:\\
		Razumevanje pojavov, procesov, motivacij, stališč in konteksta.
		\item Primeri:\\
		Intervjuji (strukturirani, polstrukturirani, odprti), Fokusne skupine, Analiza dokumentov, Opazovanje, \dots
		\item Uporaba:\\
		Ko želiš poglobljeno razložiti zakaj in kako se nekaj dogaja, ne pa meriti količine.
	\end{itemize}
	\item \textbf{kvantitativne metode}
	\begin{itemize}
		\item Kaj so?\\
		Metode, ki temeljijo na številskih podatkih in statistični obdelavi.
		\item Namen:\\
		Merjenje, primerjanje, simulacije, testiranje hipotez, iskanje korelacij.
		\item Primeri:\\
		Numerične simulacije, Anketni vprašalniki z numeričnimi odgovori, Eksperimenti z merljivimi spremenljivkami, Statistična analiza (regresija, korelacija, t-test), \dots
		\item Uporaba:\\
		Ko želiš ugotoviti koliko, kako pogosto, kakšna je povezava med spremenljivkami.
	\end{itemize}
	\item \textbf{descriptivna metoda}
	 \begin{itemize}
		\item Kaj je?\\
		Metoda, ki se osredotoča na opis stanja, pojavov ali procesov, brez nujnega iskanja vzročnih povezav.
		\item Namen:\\
		Predstaviti dejstva, značilnosti, strukturo ali potek dogajanja.
		\item Primeri:\\
		Opis obstoječih podatkov (npr. statistike iz poročil), Analiza literature, Opis postopkov ali tehnologij, \dots
		\item Uporaba:\\
		Ko želiš bralcu podati jasen pregled obstoječega stanja ali sistema, brez globlje interpretacije.
	\end{itemize}
\end{itemize}

\textbf{Glavna razlika in primerjava metod dela}
\begin{itemize}
	\item \textit{Kvalitativno}: poglobljeno razumevanje, neštevilski podatki, interpretacija.
	\item \textit{Kvantitativno}: merjenje, številčni podatki, statistika.
	\item \textit{Deskriptivno}: opis, pregled, brez analize vzročnosti.
\end{itemize}


% *********************
% *** Novo poglavje ***
% *********************
\newpage
\section{Predvidena struktura naloge}

Tukaj se pripravi predvideno kazalo naloge, recimo nekako takole kakor je spodaj, seveda morate naredit po vaši predvideni nalogi

\begin{enumerate}
	\item Uvod
	\begin{enumerate}[label*=\arabic*.]
		\item Pregled področja dela
		\item Pregled literature
		\item \dots
	\end{enumerate}
	\item Materiali in metode
		\begin{enumerate}[label*=\arabic*.]
			\item Opis modela
			\item Opis naprav
			\item \dots
		\end{enumerate}
	\item Rezultati
		\begin{enumerate}[label*=\arabic*.]
			\item Rezultati simulacije
			\item Rezultati meritve
			\item Primerjava
			\item \dots
		\end{enumerate}
	\item Diskusija in zaključek
\end{enumerate}


% ******************
% *** Literatura ***
% ******************
\cleardoublepage
\singlespacing
\addcontentsline{toc}{section}{Literatura}
\printbibliography


\end{document}
